\documentclass[11pt]{article}
\usepackage[T1]{fontenc}
\usepackage{lmodern,amsmath,amsthm,amssymb,booktabs}
\usepackage[letterpaper,textwidth=6.22in,textheight=8.75in,centering]{geometry}
\usepackage[colorlinks,linkcolor=blue!45!black,citecolor=blue!45!black,urlcolor=blue!45!black]{hyperref}
\usepackage{xcolor}
\hypersetup{pdftitle={The Alon-Jaeger-Tarsi conjecture over the five-element field in dimensions ten and eleven},pdfauthor={SII MathRSI Team}}
\newtheorem{theorem}{Theorem}[section]
\newtheorem{lemma}[theorem]{Lemma}
\newtheorem{proposition}[theorem]{Proposition}
\newtheorem{corollary}[theorem]{Corollary}
\theoremstyle{definition}\newtheorem{definition}[theorem]{Definition}
\theoremstyle{remark}\newtheorem{remark}[theorem]{Remark}
\newcommand{\F}{\mathbb F_5}
\DeclareMathOperator{\wt}{wt}
\DeclareMathOperator{\GL}{GL}
\title{The Alon--Jaeger--Tarsi conjecture over the five-element field in dimensions ten and eleven}
\author{SII MathRSI Team}
\date{November 27, 2026}
\begin{document}
\maketitle
\begin{abstract}
We prove that every invertible matrix of order ten or eleven over the field of five elements sends some nowhere-zero vector to a nowhere-zero vector. In dimension ten we count the words in the graph of the matrix according to the Hamming weights of their two blocks. An integer inequality for these counts forces a word of full weight in both blocks. In dimension eleven the corresponding inequality forces a hypothetical counterexample to have a dual word supported on three coordinates in one block and two in the other. Marking those five coordinates makes their linear relation visible to a second weight enumerator. A local Fourier transform and an integer certificate then exclude the counterexample. Both arguments are formalized in Lean. Together with the known permanent-rank bound, they give the conjecture through dimension eleven for every prime field of order at least five.
\end{abstract}

\section{Introduction}
Given two bases of a finite-dimensional vector space, can one choose a linear functional that vanishes on no vector of either basis? Over an infinite field the answer follows at once: finitely many proper hyperplanes cannot cover the dual space. Over a finite field these hyperplanes can cover it. The Alon--Jaeger--Tarsi conjecture says that, for two bases, a field of four or more elements already suffices \cite[Conjecture~1]{alontarsi}. Alon and Tarsi credit Jaeger with posing the question over the five-element field \cite[p.~393]{alontarsi}. We prove two cases over this field.

\begin{theorem}\label{thm:main}
Let $n=10$ or $n=11$. For every $A\in\GL_n(\F)$ there exists $x\in(\F^\times)^n$ such that $Ax\in(\F^\times)^n$.
\end{theorem}

The two formulations agree. In coordinates given by the first basis, write the second basis as the rows of $A$. The values of a functional on the two bases then form the vectors $x$ and $Ax$. The choice of field matters even in dimension two. Over $\mathbb F_3$, if $a,b$ are both nonzero, either $a=b$ or $a=-b$. Thus the invertible map $(a,b)\mapsto(a+b,a-b)$ has no nowhere-zero point with nowhere-zero image.

Alon and Tarsi proved the conjecture in all dimensions over proper extension fields $\mathbb F_{p^r}$, $r\ge2$ \cite[Theorem~1]{alontarsi}. Nagy and Pach proved even the affine statement, which permits one prescribed forbidden value at each input and output coordinate, in every dimension for primes $p>61$ with $p\ne79$ \cite[Theorem~A in the arXiv version]{nagypach}. For smaller primes, permanent rank gives a bound on the dimension. For every prime $p\ge5$, Yu proved the strict bound $n<2^{p-2}$ \cite{yu}. This gives seven over $\F$. A preprint of Zeng, Liu and Ratnavelu includes the endpoint $n=2^{p-2}$ for primes $p\ge5$, giving eight over $\F$ \cite[Theorem~1.1 and Corollary~1.2]{zeng}. Luo et al.\ report a Lean proof through nine for every prime $p\ge5$ in the TCSExam paper \cite[Appendix~F.3, Eq.~(9)]{tcsexam}. For $p\ge7$, Yu's bound already includes dimensions ten and eleven. The advance here is for $p=5$.

To see what changes between the two dimensions, regard the pairs $(x,Ax)$ as codewords. A table that records only the two Hamming weights still detects the desired pair: it is enough to show that cell $(n,n)$ is occupied. The dimension-nine proof reported in TCSExam also uses Fourier-related weight counts and an exact certificate, together with constraints on pairs of columns \cite[Appendix~F.3]{tcsexam}. We use induction on the dimension to rule out words with small supports in the graph and its dual. In dimension ten these exclusions and a split-weight MacWilliams identity force the full-weight cell to be occupied.

The same argument in dimension eleven produces a useful alternative. If the full-weight cell were empty, the dual would contain a word with just five nonzero entries, divided $3+2$ between the blocks after possibly exchanging them. The support of this word singles out five positions. Every word of the graph obeys the same equation on them. After rescaling and permuting coordinates we can write it as
\[
       x_1+x_2+x_3+(Ax)_1+(Ax)_2=0.
\]
Two block weights tell us how many coordinates are nonzero, but do not record their values at these positions. We keep the five entries up to their natural symmetries and count the other coordinates by weight. On the graph side the marked entries sum to zero; on the dual side we can translate all five by a common scalar. Together these restrictions give a second, finite contradiction.

\section{Which small weights can occur?}
\label{sec:support}
Write $T_n=(\F^\times)^n$, and denote the number of nonzero coordinates of a vector $z$ by $\wt(z)$. The graph
\[
 C_A=\{(x,Ax):x\in\F^n\}
 \quad\text{has orthogonal complement}\quad
 C_A^\perp=\{(-A^{\mathsf T}y,y):y\in\F^n\}.
\]
Both coordinate projections on each code are bijective. We call $(\wt(x),\wt(z))$ the split weight of $(x,z)$. We call $x$ a solution for $A$ if $x$ and $Ax$ both lie in $T_n$. Such a solution is precisely the input of a word of $C_A$ of split weight $(n,n)$.

\begin{lemma}[Support exclusions]\label{lem:support}
Suppose that every invertible matrix of dimension less than $n$ over $\F$ has a solution. If $A\in\GL_n(\F)$ has none, then
\begin{enumerate}
\item $C_A$ contains no word of split weight $(i,j)$ with $1\le i,j\le3$.
\item $C_A^\perp$ contains no word of split weight $(1,1),(1,2),(2,1),(1,3),(3,1)$ or $(2,2)$.
\end{enumerate}
\end{lemma}

\begin{proof}
For the first assertion, regard the standard coordinate vectors and the rows of $A$ as two bases of a coefficient space $E$. A functional nonzero on both bases gives the required point by taking its values on the first basis. In any space of dimension less than $n$, the induction hypothesis lets us find such a functional even if we prescribe only two independent families: extend each family to a basis and use the change-of-basis matrix.

Suppose a nonzero functional $d$ has nonzero values on precisely $i$ members of the first basis and $j$ members of the second. These are the split weights of the word defined by $d$. Call these basis vectors moving and the others fixed. If $i+j\le4$, apply induction inside $\ker d$ to the two fixed families, then extend the resulting functional $L$ to $E$. Each moving vector $v$ forbids just one value of $c$ for $L+cd$, namely $-L(v)/d(v)$. At most four values are forbidden in a field of size five.

It remains to consider $i,j\in\{2,3\}$ with $i+j\ge5$. Let $U,V$ be the respective spans of the fixed vectors, of dimensions $n-i,n-j$. Let $P$ be the set of normalized moving vectors $v/d(v)$ from either basis. Then $U+\langle P\rangle=V+\langle P\rangle=E$. We shall identify at least $k=\min(i,j)-1$ pairs successively, without introducing a dependence in either fixed family. Starting with $S=0$, suppose $\dim S=r<k$ and $S$ meets each of $U,V$ trivially. The images of $P$ in $E/(U+S)$ cannot all be equal: otherwise $E$ would be spanned by $U+S$ and one more vector, although $n-i+r+1<n$. The same holds in $E/(V+S)$. Two nonconstant maps on a common set have a pair on which both change. Indeed, if the second map takes equal values on every pair separated by the first, comparison with a point in a different fiber shows that the second is constant. Adjoin to $S$ the difference of such a pair of points of $P$. Its two images are nonzero, so the new subspace is still disjoint from each fixed span. This step strictly reduces the number of distinct images of $P$ modulo $S$.

After $k$ steps the two fixed families remain independent in $E/S$. Choose a functional there by induction and pull it back to $L$ on $E$. Now at most $|P|-k\le i+j-k=\max(i,j)+1\le4$ distinct values of $-L(v/d(v))$ occur. Again some $L+cd$ avoids every zero. This proves the first assertion, including the pair $(3,3)$ used in dimension eleven.

For the second assertion, a dual word with right weight one and left weight at most three identifies a row $R$ supported on at most three coordinates. Write $w=A^{-1}e_s$, where $s$ is its row index. Delete output $s$ after restricting inputs to a hyperplane of the form $x_j=c x_k$, $c\ne0$. The induced square map is invertible whenever $w_j\ne c w_k$: a kernel vector would map to a multiple of $e_s$ and hence would be a multiple of $w$. Any nowhere-zero solution for the smaller map has all original input coordinates nonzero. It therefore suffices to choose the hyperplane so row $R$ is nonzero there on every such input.

For a row on at most two sites, take $j,k$ to include its support with $R_j\ne0$ (the one-dimensional case is immediate). We require $cR_j+R_k\ne0$ and $w_j\ne c w_k$. The latter is not the zero condition identically, since $Rw=1$. At most two of the four nonzero scalars are excluded. For a row with three terms, at least one pair $j,k$ has $R_j w_j+R_k w_k\ne0$, since the sum over all three pairs equals $2Rw=2$. Set $c=-R_k/R_j$. This hyperplane misses $w$ and cancels exactly the chosen two terms, leaving a nonzero multiple of the third coordinate. Applying the same argument to $A^{-1}$ exchanges the two blocks and covers the reverse weight pairs.

Finally, a dual word $(u,v)$ of split weight $(2,2)$ gives $u=-A^{\mathsf T}v$. The map $A$ sends $\ker u^{\mathsf T}$ isomorphically onto $\ker v^{\mathsf T}$. On each hyperplane two coordinates are nonzero scalar multiples of one another. Deleting one of them on each side identifies the restriction with an invertible $(n-1)$-dimensional map; reinserting it preserves nonvanishing on both sides. This proves the second assertion.
\end{proof}

\section{Ordinary weights and the ten-dimensional theorem}
\label{sec:ordinary}
Let $H_{ij}$ and $D_{ij}$ count the words of split weight $(i,j)$ in $C_A$ and $C_A^\perp$, respectively. Their row and column sums are the integers $\mu_i=\binom ni4^i$, indexed by the weight in the relevant block. We use the symmetric tables $W=H+H^{\mathsf T}$ and $V=D+D^{\mathsf T}$. They have origin entry $m=2$ and row sums $m\mu_i$. In particular $W_{nn}=2H_{nn}$; this symmetrization preserves the question we want to answer.

For $0\le a,i\le n$ put
\[
 K_{ai}=[z^a](1+4z)^{n-i}(1-z)^i.
\]
If $\chi$ is a nontrivial additive character of $\F$, then for a fixed word $x$ of weight $i$ the sum of $\chi(u\cdot x)$ over $\wt(u)=a$ is $K_{ai}$. A zero coordinate of $x$ contributes four characters of value one at degree one; at a nonzero coordinate their sum is $-1$. Summing the character identity
$\sum_{c\in C_A}\chi(u\cdot c)=5^n{\bf1}_{u\in C_A^\perp}$
over an ambient split-weight class, and then adding the transposed identity, gives
\begin{equation}
 5^n V=K W K^{\mathsf T}.
 \label{eq:ordinary-transform}
\end{equation}
MacWilliams showed how the ordinary weight distribution of a binary linear code determines the distribution of its dual \cite{macwilliams}. Here the same character argument keeps the two block weights separate, which is what allows us to single out the cell $(n,n)$.

\begin{proposition}[An integer certificate]\label{prop:ordinary}
Let $W,V$ be symmetric nonnegative tables with $W_{00}=V_{00}=m>0$ and row sums $m\mu_i$ in each table. In dimensions $6\le n\le10$, \eqref{eq:ordinary-transform} and the support exclusions of Lemma~\ref{lem:support} imply $W_{nn}>0$. In dimension $11$ the same conclusion holds under the additional assumption $V_{23}=0$.
\end{proposition}
\begin{proof}
Choose the two integer vectors $u,v$ and the integer matrix $T$ specified for $n$ in the accompanying certificate. Set
\begin{align*}
 P_{ij}&=u_i+u_j+(K^{\mathsf T}TK)_{ij},&
 Q_{ij}&=v_i+v_j-5^nT_{ij},\\
 R&=2\sum_i (u_i+v_i)\mu_i .
\end{align*}
With $\langle X,Y\rangle=\sum_{ij}X_{ij}Y_{ij}$, pairing \eqref{eq:ordinary-transform} with $T$ cancels the two transform terms. The marginals give the exact identity
\[
 \langle W,P\rangle+\langle V,Q\rangle=mR.
\]
Suppose $W_{nn}=0$. Apart from the two origins, all coefficients $P_{ij},Q_{ij}$ are nonnegative on cells permitted to carry mass. In dimension eleven the test for $Q$ also exempts $(2,3),(3,2)$; they have zero mass by the extra hypothesis and symmetry. Consequently the left side is at least $m(P_{00}+Q_{00})$. Table~\ref{tab:gaps} gives the strict difference per unit mass. Each positive difference contradicts the identity.
\end{proof}

\begin{table}[ht]
\centering
\caption{Integer gaps $P_{00}+Q_{00}-R$. Only the eleven-dimensional entry requires $V_{23}=0$.}\label{tab:gaps}
\begin{tabular}{rr@{\qquad}rr}\toprule
$n$ & Gap & $n$ & Gap \\\midrule
6 & 7812500000 & 9 & 286749531250000000000\\
7 & 5860862993750000000 & 10 & 235472537548828125000000\\
8 & 1256583418475000000000 & 11 & 170253422831819243164062500000\\
\bottomrule\end{tabular}
\end{table}

For clarity, the ordinary certificate uses only a subset of Lemma~\ref{lem:support}. The primal sign test exempts the origin, $(n,n)$, and $1\le\min(i,j)\le2$, $\max(i,j)\le3$. The dual test exempts the origin and the pairs $\min(i,j)=1$, $\max(i,j)\le2$, together with the extra pair and its transpose for $n=11$. All signs and all six gaps can be checked from the integer arrays by the short replay in Appendix~\ref{sec:replay}. They require no optimization or rounding.

\begin{corollary}\label{cor:ten}
Every $A\in\GL_n(\F)$ with $n\le10$ has a solution. If an eleven-dimensional counterexample exists, its dual code contains a word with block weights $(2,3)$ or $(3,2)$.
\end{corollary}
\begin{proof}
For $n\le5$ the assertion follows from Yu's bound \cite{yu}. In a least failing dimension between six and ten, Lemma~\ref{lem:support} applies. Proposition~\ref{prop:ordinary} now gives $H_{nn}>0$, a contradiction. In dimension eleven the proposition implies $V_{23}>0$ in a counterexample. Since $V_{23}=D_{23}+D_{32}$, at least one of these counts in the actual dual code is positive.
\end{proof}

\section{A marked certificate in dimension eleven}
\label{sec:marked}

The ordinary weight distribution records the two block weights of a graph code. In dimension eleven we refine this distribution around the support of a short dual word. Orthogonality to that word restricts the five marked coordinates of every primal word. Translation by the same word imposes an additional symmetry on the dual distribution. These two facts allow a small integer certificate on the refined tables.

\subsection{Orientation and normalization}

For \(A\in\GL_{11}(\F)\), write \(C_A=\{(x,Ax):x\in\F^{11}\}\), so that \(C_A^\perp=\{(-A^{\mathsf T}y,y):y\in\F^{11}\}\). Suppose that \(A\) has no AJT vector. The ordinary certificate in Section~\ref{sec:ordinary} gives a word in \(C_A^\perp\) of block weights \((3,2)\) or \((2,3)\). In the second case exchange the two blocks. The exchanged primal code is \(C_{A^{-1}}\), and a nowhere-zero pair in it would give a nowhere-zero pair in \(C_A\). We thus work with a counterexample whose dual has a word of weights \((3,2)\).

For \(r\le 11\) let \(\mathbf 1_r\in\F^{11}\) have ones in its first \(r\) positions and zeros elsewhere. A monomial coordinate change on each of the two blocks moves the chosen dual word to \(h=(\mathbf 1_3,\mathbf 1_2)\). More explicitly, if the primal coordinate change is \((x,z)\mapsto(Lx,Rz)\), where \(L,R\) are monomial matrices, the transformed code is the graph of \(B=RAL^{-1}\) and dual words transform by \(L^{-\mathsf T}\) and \(R^{-\mathsf T}\). These latter maps can move each support to its initial positions and rescale its entries to one. The transformation preserves invertibility and the existence of a nowhere-zero pair. We therefore have a counterexample \(B\) satisfying
\begin{equation}\label{eq:marked-normalized}
 -B^{\mathsf T}\mathbf 1_2=\mathbf 1_3.
\end{equation}
We write \(C=C_B\) and mark the first three left coordinates and first two right coordinates.

\subsection{Local classes and enumerators}

Split the five marked coordinates as \(3+2\). By \eqref{eq:marked-normalized}, the marked tuple of a primal word belongs to
\[
 H=\{t\in\F^5:\textstyle\sum_{i=1}^5t_i=0\}.
\]
Let \(G\) act by permutations within each marked block and multiplication of the entire tuple by a common nonzero scalar. Classify primal marked tuples by \(G\)-orbits on \(H\). Classify dual marked tuples on all of \(\F^5\) by allowing the same actions and the shifts \(u\mapsto u+a\mathbf 1\), where \(\mathbf 1=(1,1,1,1,1)\). The permutations and scaling organize the weights and the Fourier kernel. Shifts act on the dual code itself because \(h\in C^\perp\).

There are \(29\) classes on each side. Here is a finite enumeration rule that also determines their cardinalities. For a tuple \(t\) and \(s\in\F^\times\), form \(s t+a\mathbf 1\) and sort each block separately, using the integer order on the residues \(0,\ldots,4\). Take the lexicographically least concatenation over \(s\) and the allowed \(a\), using only \(a=0\) on the primal side and all five \(a\) on the dual side. Two tuples give the same key exactly when they belong to the same class, since the sort is a choice of representative for the within-block permutations. Counting keys among the \(5^4=625\) tuples of \(H\) and the \(5^5=3125\) dual tuples gives
\begin{center}
\begin{tabular}{lrrrrrr}
\toprule
Number of classes of a given size & \(1\) & \(2\) & \(2\) & \(6\) & \(14\) & \(4\)\\
\midrule
Primal class size & \(1\) & \(4\) & \(8\) & \(12\) & \(24\) & \(48\)\\
Dual class size   & \(5\) & \(20\) & \(40\) & \(60\) & \(120\) & \(240\)\\
\bottomrule
\end{tabular}
\end{center}
For example, the dual class of \((0,0,0,0,1)\) consists of tuples with left block \((a,a,a)\) and right block \((a,a+s)\) in either order, with \(a\in\F\) and \(s\in\F^\times\). Its size is \(40\).

Use \(O=\{0,\ldots,28\}\) to index the two sets of classes separately, with zero the class of the zero tuple on each side. Let \(X=\{0,\ldots,8\}\) and \(Y=\{0,\ldots,9\}\) index the weights on the unmarked left and right coordinates, respectively. Let \(W_{obd}\) count primal words of class \(o\) and unmarked weights \(b,d\), and let \(V_{pbd}\) count dual words of dual class \(p\) with the same unmarked weights. Each table has \(29\cdot9\cdot10=2610\) cells. The corner \(W_{000}=1\) counts the zero word. The dual zero class consists of the five constant tuples, so \(V_{000}=5\), counting the multiples of \(h\). In particular, the class index \(0\) on the two sides has different meanings beyond the zero tuple itself.

\subsection{Fourier transform and marginals}

For a primal class \(o\), choose a representative \(t_o\). For a dual class \(p\), write \(\mathcal U_p\) for its set of tuples. Define
\begin{equation}\label{eq:marked-local-kernel}
 K_{po}
 =
 \#\{u\in\mathcal U_p:u\cdot t_o=0\}
 -
 \#\{u\in\mathcal U_p:u\cdot t_o=1\}.
\end{equation}
If \(\zeta\) is a primitive fifth root of unity, the sum of \(\zeta^{u\cdot t_o}\) over \(\mathcal U_p\) equals \(K_{po}\). Indeed, multiplication by a nonzero scalar permutes the class, so the four nonzero dot products all occur equally often, and \(\zeta+\zeta^2+\zeta^3+\zeta^4=-1\). Replacing \(t_o\) by another representative changes the dot products by a permutation of \(\mathcal U_p\), through the transpose of the corresponding block permutation and scalar action. For the example above, take \(t_o=(0,0,0,1,4)\). The dot products are \(\pm s\), with each nonzero value appearing ten times. Therefore this kernel entry is \(-10\).

For an unmarked block of length \(m\) define
\[
 P_m(k,b)=
 \sum_{i=0}^k(-1)^i4^{k-i}\binom bi\binom{m-b}{k-i}.
\]
For a fixed vector of weight \(b\) this is the sum of the characters paired with all target vectors of weight \(k\). At a zero coordinate the generating factor is \(1+4z\). At a nonzero coordinate it is \(1-z\), since the character sum over the four nonzero choices is \(-1\). Taking the coefficient of \(z^k\) gives the displayed formula. Fourier orthogonality for \(C\) now yields
\begin{equation}\label{eq:marked-fourier}
 5^{11}V_{pkl}
  =\sum_{o\in O}\sum_{b\in X}\sum_{d\in Y}
       W_{obd}\,K_{po}\,P_8(k,b)\,P_9(l,d).
\end{equation}
To see the equality, sum \(\sum_{c\in C}\zeta^{u\cdot c}\) over every ambient vector \(u\) of class \(p\) on the marked sites and weights \((k,l)\) off them. The inner sum is \(5^{11}\) when \(u\in C^\perp\) and zero otherwise. Summing instead in the opposite order factors the character sum into the three kernels on the right.

Write \(L_o,R_o\) for the marked left and right weights in primal class \(o\). For a dual representative \(u_p\), let \(L_p(a),R_p(a)\) be the weights of its marked blocks after adding \(a\mathbf 1\). Put \(S_m(j)=4^j\binom mj\) and define
\[
 M_L(j,b)=S_3(j)S_8(b),\qquad
 M_R(j,d)=S_2(j)S_9(d).
\]
The projections of \(C\) and \(C^\perp\) onto either block are bijections. On the primal side this gives
\begin{align}
 \sum_{o,d} [L_o=j]\, W_{obd}&=M_L(j,b),&
 \sum_{o,b} [R_o=j]\, W_{obd}&=M_R(j,d).
 \label{eq:marked-primal-marg}
\end{align}
Here \([E]\) is \(1\) when \(E\) holds and \(0\) otherwise. On the dual side put
\(
 N_L(p,j)=\#\{a:L_p(a)=j\}
\)
and similarly \(N_R(p,j)=\#\{a:R_p(a)=j\}\).
These values are invariant on a dual class: an affine scalar change on the marked tuple permutes the five shifts, and the remaining block permutations preserve each weight. The marginal identities are
\begin{align}
 \sum_{p,d}N_L(p,j)V_{pbd}&=5M_L(j,b),&
 \sum_{p,b}N_R(p,j)V_{pbd}&=5M_R(j,d).
 \label{eq:marked-dual-marg}
\end{align}
Indeed, for each shift \(a\), addition of \(ah\) permutes the dual code and leaves both unmarked weights fixed. Summing these five copies of the corresponding projection count gives \eqref{eq:marked-dual-marg}.

\subsection{Forbidden cells and the integer contradiction}

On the primal side a cell vanishes if the total weights \((L_o+b,R_o+d)\) lie on either axis away from \((0,0)\), equal \((11,11)\), or belong to
\begin{equation}\label{eq:marked-primal-small}
 \{(i,j):i,j>0,\ i+j<5\}\ \cup\ \{(2,3),(3,2),(3,3)\}.
\end{equation}
The axes use invertibility, the cell of total weights \((11,11)\) uses the counterexample hypothesis, and \eqref{eq:marked-primal-small} follows from the smaller-dimensional obstructions of Lemma~\ref{lem:support}. On the dual side a total-weight pair is forbidden if it lies on either axis away from \((0,0)\), or lies in
\[
 \{(1,1),(1,2),(2,1),(1,3),(3,1),(2,2)\}.
\]
The latter exclusions likewise use Lemma~\ref{lem:support}. A dual cell is therefore zero if \((L_p(a)+b,R_p(a)+d)\) is forbidden for any \(a\in\F\). To check this implication for an arbitrary representative of \(p\), use the fact that affine scalar changes permute the five shifted weight pairs. A word in the cell would then have a translate by \(ah\) of forbidden total weights.

We use integer multipliers \(\phi_{pkl}\) for \eqref{eq:marked-fourier}, \(w_L,w_R,v_L,v_R\) for the four marginals, and \(w_c,v_c\) for the corner equations. The full arrays are fixed in \texttt{PolynomialMarkedValues.lean}. Their entries combine into the following coefficients, where \(\delta\) is the indicator of the corner:
\begin{align}
 \alpha_{obd}
 &=\sum_{p,k,l} K_{po}P_8(k,b)P_9(l,d)\phi_{pkl}
   +w_L(L_o,b)+w_R(R_o,d)+w_c\delta_{obd}, \notag\\
 \beta_{pbd}
 &=-5^{11}\phi_{pbd}
   +\sum_{a\in\F}\bigl(v_L(L_p(a),b)+v_R(R_p(a),d)\bigr)
   +v_c\delta_{pbd}.
 \label{eq:marked-multipliers}
\end{align}
By \eqref{eq:marked-fourier} the terms in \(\phi\) cancel after pairing against \(W,V\). Equations \eqref{eq:marked-primal-marg} and \eqref{eq:marked-dual-marg} and the corners give
\begin{equation}\label{eq:marked-rhs}
 \sum_{o,b,d}W_{obd}\alpha_{obd}+\sum_{p,b,d}V_{pbd}\beta_{pbd}
 = w_c+5v_c+\langle w_L,M_L\rangle+\langle w_R,M_R\rangle
   +5\langle v_L,M_L\rangle+5\langle v_R,M_R\rangle ,
\end{equation}
with \(\langle\cdot,\cdot\rangle\) the entrywise pairing. For the prescribed arrays the corner terms \(w_c=10^{31}\) and \(5v_c=-10^{31}\) cancel. Integer evaluation of the four marginal terms gives
\begin{center}
\begin{tabular}{lr}
\toprule
Term & Value \\
\midrule
\(\langle w_L,M_L\rangle\) & \(3443757053923174209775368141760\)\\
\(\langle w_R,M_R\rangle\) & \(-10042287594608000000000056523360\)\\
\(5\langle v_L,M_L\rangle\)& \(8125598499961871048877267250400\)\\
\(5\langle v_R,M_R\rangle\)& \(-1758096222910000000000167210200\)\\
\midrule
Right-hand side of \eqref{eq:marked-rhs} & \(-231028263632954741347588341400\)\\
\bottomrule
\end{tabular}
\end{center}

The same certificate satisfies the required signs on every cell that may remain nonzero:
\begin{center}
\begin{tabular}{lrrrl}
\toprule
Table & Forbidden & Allowed & Zeros on allowed cells & Least positive coefficient \\
\midrule
Primal & \(168\) & \(2442\) & \(1\) & \(39039288019293480\)\\
Dual & \(233\) & \(2377\) & \(0\) & \(1708473306757384580\)\\
\bottomrule
\end{tabular}
\end{center}
These counts and values were independently replayed by decoding the integer arrays in the formal sources. For fixed \(p,k\), the source stores the row of \(\phi_{pkl}\) as one nonnegative integer. Extracting the \(l\)-th base-\(2^{77}\) digit and subtracting \(2^{76}\) recovers the signed entry. The two adjoint arrays use the same representation with widths \(84\) and \(95\), respectively. The replay enumerated all \(625+3125\) local tuples, recalculated the \(29^2=841\) entries of \(K\) from \eqref{eq:marked-local-kernel}, and checked both intermediate adjoint tables of \(2610\) cells apiece. The adjoints first contract against \(K\), then against \(P_9\); the final contraction against \(P_8\) gives the transform term in \(\alpha\). The sign test evaluated every cell on both sides, with no sign condition needed on forbidden cells. This replay checks the integer arithmetic. The Lean development also verifies that the structural and Fourier equations apply to the graph of an arbitrary hypothetical counterexample.

Every summand on the left of \eqref{eq:marked-rhs} is nonnegative: the forbidden cells vanish, and the remaining cells and their coefficients are nonnegative. The negative right-hand side is a contradiction. Hence a counterexample in dimension eleven over \(\F\) cannot exist.


\begin{proof}[Proof of Theorem~\ref{thm:main}]
Corollary~\ref{cor:ten} settles dimension ten. For an eleven-dimensional counterexample the same corollary supplies a dual word of the required shape. Section~\ref{sec:marked} shows that its marked counts would have a negative weighted sum in which every term is nonnegative. This rules out the remaining dimension.
\end{proof}

\section{Consequences and the next dimension}
Combining Theorem~\ref{thm:main} and Corollary~\ref{cor:ten} with Yu's theorem gives a single statement valid for every prime $p\ge5$: every invertible matrix over $\mathbb F_p$ of order at most eleven has a nowhere-zero point with nowhere-zero image. Allowing independent forbidden values in the two blocks costs no dimension.

\begin{corollary}\label{cor:affine}
Let $p\ge5$ be prime, $n\le11$, and $A\in\GL_n(\mathbb F_p)$. For any $a,b\in\mathbb F_p^n$ there is an $x$ with $x_i\ne a_i$ and $(Ax)_i\ne b_i$ for every $i$.
\end{corollary}
\begin{proof}
Write $x=a+y$ and $t=Aa-b$. We need both $y$ and $Ay+t$ nowhere zero. When $n>0$, fix an index $i$. The determinant of $B=A+c\,t e_i^{\mathsf T}$ is affine linear in $c$ and nonzero at $c=0$. At most one value of $c$ makes $B$ singular, so some nonzero $c$ remains. The homogeneous result supplies $z$ with $z$ and $Bz$ nowhere zero. Set $y=z/(c z_i)$. Then $Ay+t=Bz/(c z_i)$, as required. The case $n=0$ is immediate.
\end{proof}

The accompanying formalization also gives the affine conclusion for symmetric invertible matrices through dimension $23$. Fix a prime $p\ge5$ and suppose the affine property holds for every invertible matrix over $\mathbb F_p$ through dimension $k$. The transpose argument in the supplement then gives the following conclusion for any invertible $A$ of dimension at most $2k+1$: one can choose $A$ or $A^{\mathsf T}$ once, and that choice works for all forbidden vectors $a,b$. Taking $k=11$ gives the assertion for symmetric matrices, for which the two orientations agree. Thus at dimension twelve the remaining question over $\F$ concerns general nonsymmetric matrices.

A difficulty with going one dimension further is already visible in the ordinary counts. In dimension twelve there is a symmetric nonnegative rational table $H$ with the true integer marginals $\mu_i$, the correct origin and axes, and $H_{12,12}=0$. In fact its positive-index cells vanish if $i+j\le5$ or if $\min(i,j)=1$ and $\max(i,j)\le5$. Its formal transform $5^{-12}KHK^{\mathsf T}$ is nonnegative, has the same marginals and axes, and vanishes when $1\le i,j\le2$. The accompanying rational data give an exact example of this system of constraints. Its entries are formal multiplicities; there is no matrix construction behind the example. To obtain a twelve-dimensional proof by the marked method, one possible starting point is an argument forcing a small dual support in every counterexample, as Corollary~\ref{cor:ten} does in eleven dimensions. We do not know such an argument in dimension twelve.

\appendix
\section{Checking the certificates}\label{sec:replay}
The \href{https://github.com/lisresearchgroup/math-rsi/pull/67}{electronic supplement and public discussion} accompany this paper. The supplement contains the integer multipliers, their finite verification, and the connection to the coordinate statement of Theorem~\ref{thm:main}. We indicate how to repeat each part. All source paths below are relative to the directory containing this paper.

\paragraph{Ordinary certificates.} Run \texttt{python3 replay\_ordinary.py}. This standard-library script reads $u,v,T$ from the Lean source arrays, computes $K$ from its polynomial coefficients and checks both sign conditions in every cell before recomputing the six gaps in Table~\ref{tab:gaps}. The ordinary multipliers are in two files under \texttt{formal/Task/}: \texttt{Polynomial\-Certificate\-Values.lean} for six through ten, and \texttt{Polynomial\-Certificate\-Eleven.lean} for eleven.

\paragraph{Five marked coordinates.} Run \texttt{python3 replay\_marked.py} for an independent integer check of the local classes, Fourier kernel and all marked coefficients. The local class assignments, dot-product counts and integer multipliers are in the \texttt{PolynomialMarked} files under \texttt{formal/Task/}. The row and column order used in those arrays is fixed by the supplied class assignments. The finite sign tests are in \texttt{PolynomialMarkedValues.lean}. The link from an arbitrary counterexample to those particular tables is proved in \texttt{MarkedCounterexample.lean} and \texttt{DimensionEleven.lean}.

\paragraph{Coordinate statement.} In \texttt{formal/}, run \texttt{lake exe cache get} to fetch the pinned Mathlib build cache if desired, then \texttt{bash check.sh}. The latter builds the bundled project and checks two independent callers against the original coordinate conditions. The supplied build record uses Lean 4.33.0-rc1. All 44 local dependency modules and the root module were rebuilt from their delivered sources in an empty project build directory, reusing a pinned Mathlib cache. The nineteen printed declarations have just the standard Lean axioms \texttt{propext}, \texttt{Classical.choice} and \texttt{Quot.sound}. For the rational example in dimension twelve, run \texttt{python3 relaxation/verify.py}; that calculation uses integer arithmetic after clearing denominators.

\begin{thebibliography}{9}

\bibitem{alontarsi}
N.~Alon and M.~Tarsi.
\newblock A nowhere-zero point in linear mappings.
\newblock \emph{Combinatorica}, 9(4):393--395, 1989.
\newblock \href{https://doi.org/10.1007/BF02125351}{doi:10.1007/BF02125351}.

\bibitem{tcsexam}
Y.~Luo, Z.~Liu, Y.~Yuan, Z.~Yang, and A.~C.-C.~Yao.
\newblock TCSExam: Evaluating frontier LLMs on the full range of TCS research.
\newblock \href{https://arxiv.org/abs/2611.07962v1}{arXiv:2611.07962v1}, 5~November 2026.

\bibitem{macwilliams}
J.~MacWilliams.
\newblock A theorem on the distribution of weights in a systematic code.
\newblock \emph{Bell System Technical Journal}, 42(1):79--94, 1963.
\newblock \href{https://doi.org/10.1002/j.1538-7305.1963.tb04003.x}{doi:10.1002/j.1538-7305.1963.tb04003.x}.

\bibitem{nagypach}
J.~Nagy and P.~P.~Pach.
\newblock The Alon--Jaeger--Tarsi conjecture via group ring identities.
\newblock \emph{Journal of the European Mathematical Society}, published online first, 2025.
\newblock \href{https://doi.org/10.4171/JEMS/1640}{doi:10.4171/JEMS/1640}. \href{https://arxiv.org/abs/2107.03956v1}{Author version, arXiv:2107.03956v1}.

\bibitem{yu}
Y.~Yu.
\newblock The permanent rank of a matrix.
\newblock \emph{Journal of Combinatorial Theory, Series A}, 85(2):237--242, 1999.
\newblock \href{https://doi.org/10.1006/jcta.1998.2904}{doi:10.1006/jcta.1998.2904}.

\bibitem{zeng}
Z.~Zeng, H.~Liu, and K.~Ratnavelu.
\newblock Closing the permanent-rank boundary for the Alon--Jaeger--Tarsi conjecture.
\newblock Preprints.org, version~1, 18~August 2026.
\newblock \href{https://doi.org/10.20944/preprints202608.1194.v1}{doi:10.20944/preprints202608.1194.v1}.

\end{thebibliography}

\end{document}
